\documentclass[border=2pt,svgnames,tikz]{standalone}
\usepackage{fontspec}
\setmainfont{Source Serif 4}
\setsansfont{Source Sans 3}
\setmonofont{Source Code Pro}
\usepackage{ctex}
\setCJKmainfont[BoldFont=Source Han Sans CN]{Source Han Serif CN}
\usetikzlibrary{positioning, calc, arrows.meta, shapes.geometric}

\begin{document}
\begin{tikzpicture}[
    node distance=0.5cm and 0.5cm,
    sector/.style={draw, rectangle, minimum width=1cm, minimum height=1cm, align=center, font=\sffamily\bfseries, fill=white},
    bad_sector/.style={sector, fill=LightSalmon, draw=Red, label={[text=Red, font=\scriptsize]below:Bad}},
    spare_sector/.style={sector, fill=LightGreen, draw=Green, label={[text=Green, font=\scriptsize]below:Spare}},
    logical_label/.style={font=\small\sffamily, color=DimGray},
    physical_label/.style={font=\small\sffamily, color=DimGray},
    mapping_line/.style={->, thick, >=Stealth, color=RoyalBlue}
]

% --- Scenario 1: Sector Sparing (备用扇区/重映射) ---
\node[font=\large\bfseries\sffamily] (title1) at (0,0) {扇区备用 (Sector Sparing)};

% Logical Sectors
\node[sector] (l1_0) at (-3, -1.5) {0};
\node[sector, right=0cm of l1_0] (l1_1) {1};
\node[sector, right=0cm of l1_1] (l1_2) {2};
\node[sector, right=0cm of l1_2] (l1_3) {3};
\node[left=0.5cm of l1_0, logical_label] {Logical:};

% Physical Sectors
\node[sector, below=1.0cm of l1_0] (p1_0) {0};
\node[sector, right=0cm of p1_0] (p1_1) {1};
\node[bad_sector, right=0cm of p1_1] (p1_bad) {X};
\node[sector, right=0cm of p1_bad] (p1_3) {3};
\node[spare_sector, right=1cm of p1_3] (p1_spare) {S};
\node[left=0.5cm of p1_0, physical_label] {Physical:};

% Mappings
\draw[mapping_line] (l1_0) -- (p1_0);
\draw[mapping_line] (l1_1) -- (p1_1);
\draw[mapping_line] (l1_2.south) -- ++(0,-0.3) -| (p1_spare.north);
\draw[mapping_line] (l1_3) -- (p1_3);

% --- Scenario 2: Sector Slipping (扇区滑动) ---
\node[font=\large\bfseries\sffamily, below=4.5cm of title1] (title2) {扇区滑动 (Sector Slipping)};

% Logical Sectors
\node[sector] (l2_0) at (-3, -6.5) {0};
\node[sector, right=0cm of l2_0] (l2_1) {1};
\node[sector, right=0cm of l2_1] (l2_2) {2};
\node[sector, right=0cm of l2_2] (l2_3) {3};
\node[left=0.5cm of l2_0, logical_label] {Logical:};

% Physical Sectors
\node[sector, below=1.0cm of l2_0] (p2_0) {0};
\node[sector, right=0cm of p2_0] (p2_1) {1};
\node[bad_sector, right=0cm of p2_1] (p2_bad) {X};
\node[sector, right=0cm of p2_bad] (p2_2) {2};
\node[sector, right=0cm of p2_2] (p2_3) {3};
\node[left=0.5cm of p2_0, physical_label] {Physical:};

% Mappings
\draw[mapping_line] (l2_0) -- (p2_0);
\draw[mapping_line] (l2_1) -- (p2_1);
\draw[mapping_line] (l2_2) -- (p2_2);
\draw[mapping_line] (l2_3) -- (p2_3);

% Annotations
\node[right=0.5cm of p1_spare, font=\small\sffamily, align=left] {
    \textbf{重映射}: \\
    逻辑扇区 2 指向备用区 \\
    (通常用于磁盘使用中产生的坏块)
};

\node[right=0.5cm of p2_3, font=\small\sffamily, align=left] {
    \textbf{滑动}: \\
    坏块后的所有扇区物理位置后移 \\
    (通常用于低级格式化时)
};

\end{tikzpicture}
\end{document}
